%----------------------------------------------------------------------------- %%
\chapter{\babSatu}
\label{bab:1}
%----------------------------------------------------------------------------- %%
This chapter provides the main background and challenges that drive this thesis. Section 1.1 outlines the research background, highlighting the motivation behind addressing resource fragmentation and network locality in Kubernetes clusters. Section 1.2 defines the problem statement, specifying the main research questions and limitations. Section 1.3 enumerates the objectives of the thesis, while Section 1.4 describes the anticipated research contributions. Section 1.5 positions this work within relevant academic domains. Section 1.6 summarizes the principal steps taken during the research. This chapter concludes with Section 1.7, which presents the thesis outline.

\section{Research Background}
\label{sec:researchBackground}

Kubernetes schedules Pods by filtering feasible Nodes and scoring the remaining alternatives. This process is effective for initial placement, but running Pods normally remain bound to their Nodes until they terminate or are recreated. As workloads scale and change, free CPU and memory can become scattered into incompatible shapes. Aggregate capacity may be sufficient while no individual Node can accommodate a new Pod, a condition commonly described as resource fragmentation or the ``Tetris problem''~\citep{mishra_bellur_2016,grandl_tetris_2014}.

Fragmentation affects schedulability, utilization, and infrastructure cost. A Node may retain CPU without enough memory, while another retains memory without enough CPU. These asymmetric residuals strand capacity because the two dimensions cannot be allocated independently~\citep{xiao2013dynamic,guo2025joint}. Directed Pod eviction can repair such placement states and reduce the number of Nodes required by the same workload~\citep{larsson_klein_elmroth_2023}.

The Kubernetes Descheduler provides a post-placement correction mechanism by evicting selected Pods and allowing the kube-scheduler to place their replacements~\citep{k8s_descheduler_docs}. For node reclamation, the important question is not merely whether a source Node is lightly utilized, but whether an evicted Pod has a feasible target that produces denser and better-shaped packing. This requires source-node assessment, target prediction, and an eviction budget that limits disruption.

Resource packing can also affect application communication. Moving dependent microservices to distant or degraded Nodes may improve compute packing while increasing latency. Communication-aware placement studies show that network topology, observed latency, and service dependencies can materially affect placement quality~\citep{marchese2022network,marchese_tomarchio_2022,chen2024trade}. A practical node reclamation architecture therefore also needs a network-aware guardrail.

This research develops a defragmentation-focused Descheduler architecture with two complementary components. \textit{ResourceDefragmentation} identifies drain candidates using utilization and CPU-memory shape, predicts feasible scheduler targets, and chooses a Pod using the projected bin score of its predicted landing Node. A network-aware filtering layer blocks evictions whose feasible destinations would degrade communication locality.

\section{Problem Statement}
\label{sec:problemStatement}

The study addresses the following problems:

\begin{enumerate}
    \item \textbf{Scattered capacity:} CPU and memory remain available in aggregate but cannot satisfy workload requests on any single Node.
    \item \textbf{Stranded resource dimensions:} CPU-heavy and memory-heavy placement states leave asymmetric capacity that is difficult to reuse.
    \item \textbf{Missing post-placement node reclamation:} Scheduler bin-packing scores act when a Pod is placed and do not repair an already fragmented cluster.
    \item \textbf{Uncertain rescheduling outcome:} An eviction is useful only if the Pod has a feasible non-origin target that improves packing without exceeding a utilization ceiling.
    \item \textbf{Network-locality risk:} Compute-efficient movement can separate communicating services and degrade application latency.
\end{enumerate}

\subsection{Research Question}
\label{sec:researchQuestion}

\begin{itemize}
    \item How can a Descheduler policy identify Nodes whose low utilization or poor CPU-memory shape makes them suitable drain candidates?
    \item How can predicted scheduler targets and a single target-bin-score criterion select effective and safe eviction candidates?
    \item How effectively does the proposed policy reclaim Nodes, reduce stranded capacity, preserve balanced headroom, and converge compared with \textit{HighNodeUtilization}?
    \item How can network-aware filtering reduce latency risk without preventing useful resource improvement?
\end{itemize}

\subsection{Research Limitation}
\label{sec:researchLimitation}

\begin{itemize}
    \item Evaluation uses a controlled multi-node Kubernetes deployment and synthetic workloads with declared CPU and memory requests.
    \item Resource-policy experiments use request-based accounting because kube-scheduler feasibility is request-based.
    \item Final placement remains the responsibility of the configured kube-scheduler; the Descheduler predicts but does not bind Pods.
    \item Movement uses standard Kubernetes eviction and controller recreation rather than live migration.
    \item The resource model focuses on CPU and memory; storage, accelerators, and additional extended resources are outside the main evaluation.
    \item Network awareness is implemented as eviction filtering, not as a complete routing or traffic-engineering system.
\end{itemize}

\section{Research Objectives}
\label{sec:researchObjectives}

\begin{itemize}
    \item Design a defragmentation-focused Descheduler policy that combines utilization and resource-shape assessment.
    \item Quantify dimensional skew and free-space pressure using RII, FSI, and a CPU-memory bin score.
    \item Predict feasible landing Nodes before eviction and select eviction candidates using projected target balance.
    \item Enforce an eviction budget, target-utilization ceiling, and pack-upward constraint to reduce unsafe or unproductive movement.
    \item Evaluate node reclamation, stranding reduction, balanced headroom, eviction efficiency, convergence, and safety.
    \item Integrate network-aware filtering that protects communication locality during workload movement.
\end{itemize}

\section{Research Contributions}
\label{sec:researchContributions}

\begin{itemize}
    \item \textbf{A defragmentation-focused Descheduler policy:} \textit{ResourceDefragmentation} repairs post-placement scatter without replacing the kube-scheduler.
    \item \textbf{RII- and stranding-aware Node assessment:} Drain candidates are identified from average utilization and CPU-memory bin quality, then prioritized using RII and FSI.
    \item \textbf{A predicted-target eviction-candidate selector:} Candidate Pods are compared using the projected bin score of their predicted landing Nodes.
    \item \textbf{Scheduler-target prediction and safety constraints:} Evictions require request feasibility, a target-utilization ceiling, and a target at least as packed as the source.
    \item \textbf{Network-aware eviction filtering:} Topology- or latency-based costs protect dependent services from harmful movement.
    \item \textbf{A scenario-based evaluation:} Five node reclamation and fragmentation scenarios quantify reclamation, stranding, headroom, efficiency, convergence, and safety.
\end{itemize}

\section{Research Position}
\label{sec:researchPosition}

Existing Kubernetes bin-packing strategies such as \texttt{MostAllocated} and \texttt{RequestedToCapacityRatio} improve initial placement but do not move running Pods~\citep{k8s_resource_bin_packing}. Upstream Descheduler utilization policies can pack workloads using thresholds, but a scalar utilization threshold does not directly represent CPU-memory shape. Research on directed eviction and multi-resource scheduling establishes that placement repair and dimensional balance can improve effective capacity~\citep{larsson_klein_elmroth_2023,grandl_tetris_2014,guo2025joint}.

This study positions \textit{ResourceDefragmentation} between threshold-only draining and a full scheduler replacement. The policy uses the Descheduler to choose limited, purposeful evictions while preserving kube-scheduler ownership of placement. Its defining choice is a single predicted-target criterion: among feasible Pods on a drain Node, prefer the Pod whose predicted destination has the strongest projected CPU-memory bin score. The accompanying network-aware layer treats communication cost as an eviction guardrail.

\section{Research Steps}
\label{sec:langkahPenelitian}

\begin{enumerate}
    \item Review Kubernetes scheduling, fragmentation, node reclamation, Descheduler policies, and communication-aware placement.
    \item Define Node-state metrics, drain-candidate rules, target prediction, eviction-candidate selection, and safety constraints.
    \item Implement the resource and network-aware plugins in the Kubernetes Descheduler framework.
    \item Evaluate the resource policy on five controlled scenarios and compare it with \textit{HighNodeUtilization}.
    \item Evaluate network-aware filtering using topology- and latency-sensitive workloads.
    \item Analyze reclamation, stranding, headroom, eviction efficiency, convergence, safety, and latency outcomes.
\end{enumerate}

\section{Thesis Outline}
\label{sec:sistematikaPenulisan}

\begin{enumerate}
    \item \textbf{Chapter 1} introduces the problem, questions, objectives, and contributions.
    \item \textbf{Chapter 2} reviews Kubernetes resource management, fragmentation, node reclamation, Descheduler mechanisms, and network-aware placement.
    \item \textbf{Chapter 3} presents the \textit{ResourceDefragmentation} pipeline and the Network-Aware Evictor design.
    \item \textbf{Chapter 4} evaluates node reclamation, defragmentation, safety, convergence, and network-latency trade-offs.
    \item \textbf{Chapter 5} summarizes the findings and future work.
\end{enumerate}
